\documentclass[a4paper, 12pt]{article}

\usepackage{utils}

\renewcommand*{\today}{09 March 2026}

\begin{document}

\hotbox{Algèbre bilinéaire}{CM 6}{\today}

\section{Espaces euclidiens}

\subsection{Définitions}

\begin{definition}
    Un \textbf{espace euclidien} est un espace vectoriel réel $E$ muni d'un produit scalaire.
\end{definition}

\begin{remarque}
    Dans le cas de la dimension infinie, on dit que $E$ est un \textbf{espace préhilbertien}.
\end{remarque}

\begin{definition}
    On appelle \textbf{norme euclidienne} la norme associée au produit scalaire, c'est-à-dire la fonction $\norm{x} = \sqrt{\langle x, x \rangle}$.
\end{definition}

Dans la suite, nous considérerons uniquement la norme euclidienne.

\begin{definition}
    Les \textbf{identités de polarisation} s'écrivent, pour tous $x, y \in E$,
    \begin{align*}  
        \langle x, y \rangle &= \frac{1}{2} \left( \norm{x + y}^2 - \norm{x}^2 - \norm{y}^2 \right) \\
        &= \frac{1}{2} \left( \norm{x + y}^2 - \norm{x - y}^2 \right)
    \end{align*}
\end{definition}

\begin{theoreme}{Inégalité de Cauchy-Schwarz}
    Pour tous $x, y \in E$, on a
    $$    
    |\langle x, y \rangle| \leq \norm{x} \norm{y}
    $$
\end{theoreme}

\begin{demonstration}
    Posons $P(\lambda) = \norm{x + \lambda y}^2$ pour $\lambda \in \R$.
    \begin{align*}
        P(\lambda) &= \langle x + \lambda y, x + \lambda y \rangle \\
        &= \norm{x}^2 + 2\lambda \langle x, y \rangle + \lambda^2 \norm{y}^2
    \end{align*}

    Comme $P(\lambda) \geq 0$ pour tout $\lambda$, le discriminant de ce polynôme doit être négatif ou nul, c'est-à-dire
    \begin{align*}
        4 \langle x, y \rangle^2 - 4 \norm{x}^2 \norm{y}^2 &\leq 0 \\
        \langle x, y \rangle^2 &\leq \norm{x}^2 \norm{y}^2 \\
        |\langle x, y \rangle| &\leq \norm{x} \norm{y}
    \end{align*}
\end{demonstration}

\begin{theoreme}{}{}
    Soit $x, y \in E$. On a
    \begin{enumerate}
        \item $\norm{x + y}^2 + \norm{x - y}^2 = 2(\norm{x}^2 + \norm{y}^2)$ (identité de parallélogramme)
        \item $\norm{x}^2 - \norm{y}^2 = \langle x + y, x - y \rangle$
    \end{enumerate}
\end{theoreme}

\begin{remarque}
    L'inégalité de Cauchy-Schwarz peut s'écrire $-1 \leq \dfrac{\langle x, y \rangle}{\norm{x} \norm{y}} \leq 1$.
    On peut interpréter ça comme un "angle".
\end{remarque}

\begin{remarque}[$\mathbb{C}$ c'est mieux!]
    Soit $E$ un $\mathbb{C}$-espace vectoriel muni d'une forme sesquilinéaire définit positive $\langle \cdot, \cdot \rangle$.
    (comparé à un produit scalaire, une forme sesquilinéaire a comme propriété $\langle x, \lambda y \rangle = \bar{\lambda} \langle x, y \rangle$ pour la deuxième coordonnée.)
\end{remarque}

\subsection{Procédé de Gram-Schmidt}

\begin{theoreme}{}{}
    Soit $E$ un espace euclidien de dimension finie $n$ et $(v_1, \dots, v_n)$ une base de $E$.

    Alors il existe une unique base orthonormale $(e_1, \dots, e_n)$ de $E$ telle que pour tout $1 \leq j \leq n$, on a
    \begin{enumerate}
        \item $\langle e_j, v_j \rangle > 0$
        \item $Vect(e_1, \dots, e_j) = Vect(v_1, \dots, v_j)$
    \end{enumerate}
\end{theoreme}

\begin{demonstration}
    On définit les $e_i$ récursivement de la manière suivante :
\begin{align*}
    e_1 &= \frac{v_1}{\norm{v_1}} = u_1 \\
    u_2 &= v_2 - \dfrac{\langle u_1, v_2 \rangle}{\norm{u_1}} u_1, \quad e_2 = \frac{u_2}{\norm{u_2}} \\
    u_j &= v_j - \sum_{i=1}^{j-1} \dfrac{\langle u_i, v_j \rangle}{\norm{u_i}} u_i, \quad e_j = \frac{u_j}{\norm{u_j}}
\end{align*}
\end{demonstration}

\end{document}